% Fijiisland
% 2022-06-03
\chapter{Normalization}

\mksec{命名}
为了统一文内特殊名词的使用，本译作在此页标注了一些常用的名词翻译，详见下表。

\begin{table}[h]
    \footnotesize
    \begin{tabularx}{\linewidth}{X X}
        \toprule
        \textbf{原文} & \textbf{中译}\\
        \midrule
        numeric & 数值的 \\
        stand-in notation & 代记法 \\
        \bottomrule
    \end{tabularx}
\end{table}

\mksec{排版}
为保证本译作排版井然有序，下面列出了一些写作时需遵循的标准。

\medskip
\begin{itemize}
    \item section之前使用 \redtxt{{\textbackslash}bigskip} 制作大间隔。
    \item subsection之前使用 \redtxt{{\textbackslash}medskip} 制作中间隔。
    \item 文本后的代码块不需要制作间隔。代码块之后使用 \redtxt{{\textbackslash}smallskip} 制作小间隔。
    \item figure前后使用 \redtxt{{\textbackslash}smallskip} 制作小间隔。
    \item figure的scale一般为0.29。特殊情况可酌情调整。
    \item table的间隔同figure。
\end{itemize}